%------------------------------------------%
% Cannabis Data Science
% Saturday Morning Statistics #24
% Date: 5/14/2022
%------------------------------------------%
\documentclass[xcolor={dvipsnames}]{beamer}

%------------------------------------------%
% Define the title.
%------------------------------------------%
\author{Cannabis Data Science}
\title[\textbf{Saturday Morning Statistics \#24}]{}
\institute[]{\Large Saturday Morning Statistics \#24}
\date{May \nth{14}, 2022}

%------------------------------------------%
% Define the presentation style.
%------------------------------------------%

\hypersetup{pdfpagemode = FullScreen}
\mode<presentation>{
  \usetheme{Boadilla}
  \usecolortheme{orchid}
  \usefonttheme{default}
  \setbeamertemplate{navigation symbols}{}
  \setbeamertemplate{caption}[numbered]
}
\defbeamertemplate*{title page}{customized}[1][]{
  \usebeamerfont{title}\inserttitle\par
  \bigskip
  \vspace{0.5\baselineskip}
  \usebeamerfont{institute}\insertinstitute\par
  \vspace{0.5\baselineskip}
  {\small\usebeamerfont{date}\insertdate\par}
  \usebeamercolor[fg]{titlegraphic}\inserttitlegraphic
}
\setbeamersize{
  text margin left = 0.5in,
  text margin right = 0.5in
}
%------------------------------------------%
% Specify packages.
%------------------------------------------%
\usepackage{amsmath}
\usepackage{caption}
\usepackage[english]{babel}
\usepackage{graphicx}
\usepackage{hhline}
\usepackage[utf8x]{inputenc}
\usepackage{mathtools} % Annotating equations.
\usepackage[super]{nth} % 1st, 2nd, 3rd, etc.
\usepackage{setspace}
\usepackage{subcaption}
\usepackage{tikz}
\usepackage{xparse}

%------------------------------------------%
% Set the theme.
%------------------------------------------%

% Colors
\definecolor{LG}{RGB}{218, 247, 166}
\definecolor{DG}{RGB}{2, 48, 32}

% Palette
\setbeamercolor*{palette primary}{bg=LG, fg=DG}
\setbeamercolor*{palette secondary}{bg=LG, fg=DG}
\setbeamercolor*{palette tertiary}{bg=LG, fg=DG}

%------------------------------------------%
% Define all custom commands.
%------------------------------------------%

% Top space.
\newcommand\T{\rule{0pt}{2.5ex}}

% Bottom space.
\newcommand\B{\rule[-1.25ex]{0pt}{0pt}}

% Blocks.
\newenvironment<>{Block}[2][.9\textwidth]
  {\setlength{\textwidth}{#1}
  \begin{actionenv}#3
    \def\insertblocktitle{#2}\par
    \usebeamertemplate{block begin}}
  {\par\usebeamertemplate{block end}
  \end{actionenv}}

% Balls.
\defbeamertemplate{enumerate item}{largeball}
{\begin{pgfpicture}{-1ex}{-0.65ex}{1.5ex}{1.5ex}
\usebeamercolor[fg]{item projected}
{\pgftransformscale{2.5}\pgftext{\Large\pgfuseshading{bigsphere}}}
{\pgftransformshift{\pgfpoint{0pt}{0.5pt}}
\pgftext{\usebeamerfont*{item projected}\small\insertenumlabel}}
\end{pgfpicture}}

% Fancy arrows.
\NewDocumentCommand\UpArrow{O{2.0ex} O{black}}{%
   \mathrel{\tikz[baseline] \draw [->, line width=0.5pt, #2] (0,0) -- ++(0,#1);}} % Fancy up-arrow.
\NewDocumentCommand\DownArrow{O{2.0ex} O{black}}{%
   \mathrel{\tikz[baseline] \draw [<-, line width=0.5pt, #2] (0,0) -- ++(0,#1);}} % Fancy down-arrow.

% Equations with numbers on the left.
\makeatletter
\newcommand{\LeftEqNo}{\let\veqno\@@leqno}
\makeatother

% No separating line on footnote.
\renewcommand*\footnoterule{}


%------------------------------------------%
%------------------------------------------%
% Presentation
%------------------------------------------%
%------------------------------------------%
\begin{document}

% Title page.
\begin{frame}{}

% Background.
\tikz[remember picture, overlay]
\node[opacity=1.0, inner sep=0pt] at (current page.center){
  \includegraphics[height=\paperheight, width=\paperwidth]{images/presentation-cover.pdf}
};

% Title.
\vspace*{3\baselineskip}
\includegraphics[scale=0.375]{images/logo.pdf}
\vspace*{-2\baselineskip}
\titlepage

\end{frame}


%------------------------------------------%
% Introduction
%------------------------------------------%


\begin{frame}{Enter the Skunk}

\begin{center}
\includegraphics[width=1.25in]{images/skunk.pdf}\\
\end{center}

\begin{minipage}{0.9\textwidth}
{\itshape ``"It’s been hard to breathe and the smell’s been just horrendous. [It’s] like you’ve literally been sprayed by a {\bfseries skunk}.”''}
\end{minipage}

\vspace{0.25\baselineskip}

\begin{flushright}
\begin{minipage}{0.75\textwidth}
{\footnotesize - Resident of Prague, Oklahoma\\{\itshape ``‘It’s nasty’: Prague neighbors push back on area cannabis facility}, Oklahoma News 4}\\
\end{minipage}
\end{flushright}

{\tiny\color{gray}https://kfor.com/news/local/its-nasty-prague-neighbors-push-back-on-area-cannabis-facility/}

\end{frame}

%------------------------------------------%
% Literature Review
%------------------------------------------%

\begin{frame}{Literature Review}

{\bfseries Over eight hundred cannabis strains characterized by the relationship between their psychoactive effects, perceptual profiles, and chemical compositions}

\vspace{0.5\baselineskip}

{\footnotesize by Laura Alethia de la Fuente, Federico Zamberlan, andres sanchez, Facundo Carrillo, Enzo Tagliazucchi, Carla Pallavicini}

\vspace{0.5\baselineskip}
   
\begin{itemize}

\item {\bfseries Claim}: {\itshape ``While cannabinoid content was variable even within individual strains, \underline{terpene profiles} matched the perceptual characterizations made by the users and could be used to \underline{predict} associations between different psychoactive \underline{effects}.''}


\end{itemize}

\vspace{2.5\baselineskip}

{\tiny\color{gray}https://www.biorxiv.org/content/10.1101/759696v1.abstract}

\end{frame}

%------------------------------------------%
% Question of the day
%------------------------------------------%

\section{Hypothesis}
\begin{frame}{Question and Hypothesis}

% Question of the day
\begin{center}
\begin{minipage}{.9\linewidth}
\begin{Block}{Question of the day.}

\vspace{.5\baselineskip}
\begin{itemize}

\item Can we build a model to predict a cannabis product's {\bfseries potential} \underline{effects and flavors} given the product's \underline{lab results}.

\end{itemize}

\vspace{.5\baselineskip}

\end{Block}
\end{minipage}
\end{center}


\begin{center}
\includegraphics[width=1.25in]{images/skunk.pdf}%
\hspace{0.2\textwidth}%
\includegraphics[width=1.25in]{images/lemon.pdf}
\\
\end{center}

\end{frame}



%------------------------------------------%
% Methodology
%------------------------------------------%

\begin{frame}{Methodology: Probit Models}

Given a latent variable representation of the {\bfseries probit model}:

\begin{align*}
z_i &= x_i\beta + \epsilon _i, \hspace{4ex} \epsilon _i \stackrel{iid}{\sim} \mathcal{N}(0, 1), \\
\\[-0.5\baselineskip]
y_i &=
  \begin{cases}
    1 \text{ if } z_i > 0\\
    0 \text{ if } z_i \leq 0
  \end{cases}
\end{align*}

\vspace{0.5\baselineskip}
You can estimate the parameters using the {\bfseries likelihood function}

$$
L(\beta) = \prod_{i=1}^n\Phi(x_i\beta)^{y_i}[1 - \Phi(x_i\beta)]^{1-y_i}.
$$

\end{frame}


%------------------------------------------%
% Results
%------------------------------------------%



%------------------------------------------%
% Conclusion
%------------------------------------------%


%------------------------------------------%
% Takeaway
%------------------------------------------%
\section{Takeaway}
\begin{frame}{}
\begin{center}

% Thank you.
\begin{minipage}{3.85in}
\begin{center}
\includegraphics[width=.25in]{images/prayer.png} {\Large \textbf{Thank you for coming.}}\\
\end{center}
\vspace*{0.5\baselineskip}

% Insight of the day.
\begin{center}
\begin{minipage}{\linewidth}
\begin{Block}{Insights of the Day}

\vspace{0.5\baselineskip}

\begin{itemize}

\item Stand on the shoulders of giants to reach the hard--to--get fruit.

\vspace{0.5\baselineskip}

\item Simple is better than complex.

\vspace{0.5\baselineskip}

\end{itemize}

\end{Block}
\end{minipage}
\end{center}

\vspace*{2\baselineskip}

\end{minipage}

\end{center}
\end{frame}
\end{document}
%------------------------------------------%
%------------------------------------------%